\section{The First Week of Advent: Awakening to the Coming Lord}\label{sec:week1}

The liturgical year does not open with a lullaby; it opens with an alarm bell. Across all three Sunday cycles, the First Sunday of Advent places the assembly on watch at the edge of night, summoning a drowsy world to wake from sleep, lift its head, and walk in the light of the Lord.

\subsection{The First Sunday of Advent Across Cycles A, B, and C}

\begin{threestable}{Cycle}{Appointed Readings}{Thematic Arc}
Year A & Isaiah~2:1--5; Psalm~122 (121); Romans~13:11--14a; Matthew~24:37--44 & Pilgrimage to Sion's mountain of peace; waking from sleep to put on Christ; keeping watch like the householder who knows not the thief's hour. \\
\addlinespace
Year B & Isaiah~63:16b--17, 19b; 64:2b--7 (Vulgate 63:16b--17; 64:1, 3--8); Psalm~80 (79); 1~Corinthians~1:3--9; Mark~13:33--37 & Crying out for God to rend the heavens and come down; being confirmed blameless in grace; the doorkeeper commanded to watch through every watch of the night. \\
\addlinespace
Year C & Jeremiah~33:14--16; Psalm~25 (24); 1~Thessalonians~3:12--4:2; Luke~21:25--28, 34--36 & The righteous Shoot raised up for David; abounding in charity before the Lord's coming; lifting up our heads amid cosmic shaking because redemption is near. \\
\end{threestable}

\subsubsection{Year A: Walking in the Light and Putting on Christ}

In Year~A, Isaiah opens the new year with the vision of the last days, when the mountain of the Lord's house stands above the hills, nations beat their swords into ploughshares, and the prophet calls out: ``O house of Jacob, come ye, and let us walk in the light of the Lord'' (Isaiah~2:5). Psalm~122 (Vulgate~121) answers with the pilgrim's song---``I rejoiced at the things that were said to me: We shall go into the house of the Lord''---and St.~Paul strikes the decisive waking note in Romans~13:11--14:

\begin{quote}
And that, knowing the season, that it is now the hour for us to rise from sleep. For now our salvation is nearer than when we believed. The night is passed [and] the day is at hand. Let us, therefore cast off the works of darkness and put on the armour of light. \dots\ But put ye on the Lord Jesus Christ: and make not provision for the flesh in its concupiscences.
\end{quote}

Those very sentences turned the life of St.~Augustine in a garden in Milan. Hearing a child's voice chanting ``Take up and read,'' Augustine snatched the Apostle's book and read Romans~13:13--14 in silence (\textit{Confessions} VIII.12.29): ``No further would I read, nor did I need; for instantly, as the sentence ended,---by a light, as it were, of security infused into my heart,---all the gloom of doubt vanished away.'' Preaching on the same verses in Constantinople, St.~John Chrysostom (\textit{Homily~24 on Romans}) presses the urgency home: the longer we have lived in faith, the closer stands the day of salvation, so a Christian may not dress in nightclothes when dawn is already breaking on the ridge. Matthew~24:37--44 seals the day with the Lord's sober warning: as in the days of Noah, ordinary eating, drinking, and marrying went on until the flood came; ``Watch ye therefore, because you know not what hour your Lord will come.''

\subsubsection{Year B: ``O That Thou Wouldst Rend the Heavens''}

Year~B opens with one of the rawest cries in the Old Testament. Standing amid Israel's sin and exile, Isaiah does not pretend to righteousness; he confesses that we have fallen like withered leaves and begs God to tear open the sky: ``O that thou wouldst rend the heavens, and wouldst come down: the mountains would melt away at thy presence. \dots\ And now, O Lord, thou art our father, and we are clay: and thou art our maker, and we all are the works of thy hands'' (Isaiah~64:1, 8; Lectionary~63:19b, 64:7). Psalm~80 (Vulgate~79) repeats the cry as a litany---``Stir up thy might, and come to save us. \dots\ O God of hosts, convert us: and shew thy face, and we shall be saved''---before St.~Paul reassures the Corinthians that ``God is faithful,'' who will confirm them without crime unto the day of the coming of our Lord Jesus Christ (1~Corinthians~1:8--9).

Mark~13:33--37 distills the whole Gospel of Advent into a single verb repeated like a sentinel's beat: ``Take ye heed, watch and pray. For ye know not when the time is. \dots\ Watch ye therefore (for you know not when the lord of the house cometh, at even, or at midnight, or at the cock crowing, or in the morning): Lest coming on a sudden, he find you sleeping. And what I say to you, I say to all: Watch.'' The four watches of the Roman night---evening, midnight, cockcrow, dawn---embrace every season of a human life, from youth to old age. None is exempt from the porter's post.

\subsubsection{Year C: Lifting Up Our Heads}

In Year~C, Jeremiah promises that the Lord will fulfill His good word and raise up for David ``the bud of justice'' whose name is ``The Lord our just one'' (Jeremiah~33:14--16). Psalm~25 (Vulgate~24) answers with the classic opening prayer of Roman Advent: ``To thee, O Lord, have I lifted up my soul. In thee, O my God, I put my trust; let me not be ashamed.'' St.~Paul prays that the Thessalonians may abound in charity toward one another and toward all, so that their hearts stand blameless in holiness ``at the coming of our Lord Jesus Christ, with all his saints'' (1~Thessalonians~3:12--13).

Luke~21:25--28, 34--36 brings the First Sunday's eschatological tension to a climax. Amid signs in sun, moon, and stars and the distress of nations, Jesus refuses panic and commands hope: ``But when these things begin to come to pass, look up and lift up your heads, because your redemption is at hand.'' Preaching in sixth-century Rome while plague, flood, and invasion battered the city, Pope St.~Gregory the Great (\textit{Homily~1 on the Gospels}, \S1, \S5) read this Gospel as mercy rather than cruelty: ``Our Lord and Saviour wisheth to find us ready at His second coming. Therefore He telleth us what will be the evils of the world as it groweth old, that He may wean our hearts from worldly affections.'' St.~Bede the Venerable (\textit{Exposition on Luke} 21:28--29) adds that lifting up our heads means raising our minds to the joys of our heavenly homeland: just as the fig tree's budding announces summer, so the passing of this fragile world heralds the kingdom of God, when the clouds of our sorrow pass away (\textit{moeroris nostri nubila transeunt}) and the days of life shine with the brightness of the eternal Sun. Pope St.~Leo the Great (\textit{Sermon~19.1}) applies the closing warning of Luke~21:34 directly to December: ``For the advent of which it behoves every man to prepare himself, lest it find him given over to gluttony, or entangled in cares of this life.''

\subsection{The Weekdays of the First Week}

Unlike the 1962 calendar, the postconciliar Lectionary gives every weekday of Week~I its own prophetic and Gospel pairing (\textit{Ordo Lectionum Missae}, nos.~175--180), launching a continuous reading of Isaiah's promises alongside Matthew's testimony to the Messiah's deeds:

\begin{itemize}
  \item \textbf{Monday:} Isaiah~2:1--5 (in Year~A, Isaiah~4:2--6, the ``bud of the Lord'' and the sheltering cloud and fire over Mount Sion, may be substituted because Isaiah~2 was proclaimed on Sunday) and Matthew~8:5--11, the Roman centurion's faith proving that many will come from east and west to recline at table with Abraham, Isaac, and Jacob.
  \item \textbf{Tuesday:} Isaiah~11:1--10 (the rod from the root of Jesse, the sevenfold Spirit, and the wolf dwelling with the lamb) and Luke~10:21--24, Christ rejoicing in the Holy Spirit because the Father reveals to little ones what prophets and kings longed to see.
  \item \textbf{Wednesday:} Isaiah~25:6--10a (the Lord's mountain banquet of rich food and choice wines where tears are wiped away and death is destroyed forever) and Matthew~15:29--37, Jesus healing the crowd on the mountain and feeding four thousand with seven loaves.
  \item \textbf{Thursday:} Isaiah~26:1--6 (``Open ye the gates, and let the just nation enter in,'' trusting in the Lord the everlasting rock) and Matthew~7:21, 24--27, building one's house upon rock not by saying ``Lord, Lord,'' but by doing the Father's will.
  \item \textbf{Friday:} Isaiah~29:17--24 (the deaf hearing the words of the book and the eyes of the blind seeing out of darkness) and Matthew~9:27--31, two blind men crying ``Have mercy on us, O Son of David'' and receiving sight according to their faith.
  \item \textbf{Saturday:} Isaiah~30:19--21, 23--26 (the Lord waiting to show compassion, a voice behind you saying ``This is the way, walk ye in it'') and Matthew~9:35--10:1, 5a, 6--8, Jesus moved with compassion for crowds like sheep without a shepherd and sending out the Twelve.
\end{itemize}

\subsection{Evening Prayer and the Practice of the First Week}

Every evening of the first phase of Advent, Evening Prayer in \textit{Liturgia Horarum} (1971) sings the restored seventh-century hymn \textit{Conditor alme siderum}; for household prayer, Appendix~B prints the received 1632/1888 public-domain Roman Breviary text \textit{Creator alme siderum} alongside Edward Caswall's 1849 English translation, addressing Christ as ``Creator of the starry frame, Eternal Light of all who live,'' who came from a Virgin's shrine to cure a sick world and who will return as ``Great Judge of all.'' Sing or recite one stanza beside the single lighted candle of the Advent wreath each dusk.

For the discipline of the First Week, choose a modest concrete act of watchfulness: rise fifteen minutes earlier to read the weekday Isaiah and Gospel passages before opening the news; keep Friday meatless or offer a deliberate Friday penance; and choose the person or local charity through whom you will practice a corporal work of mercy during the weeks ahead.
